\upshape
\emph{Idea.} Each threshold compares the gross gain from silence in~\eqref{eq:gains} with the hiding cost; the attention threshold rewrites $\tilde e>1$ on the cost scale. Throughout, $b>0$, as the definition $\tilde e=\Delta/b$ presupposes.

\emph{Step 1 (low-quality seller, (i)).} An informed seller with $\theta=0$ gains $g_0=b\,\tilde e$ by hiding, by~\eqref{eq:gains}, and pays $c_0$. She strictly prefers non-disclosure iff $b\,\tilde e>c_0$, that is, iff $\tilde e>c_0/b$.

\emph{Step 2 (high-quality seller, (ii)).} An informed seller with $\theta=1$ gains $g_1=b\,(\tilde e-1)$ and pays $c_0+c_1$. She strictly prefers non-disclosure iff
\begin{equation*}
    b\,(\tilde e-1)\;>\;c_0+c_1\quad\Longleftrightarrow\quad\tilde e\;>\;1+\frac{c_0+c_1}{b}.
\end{equation*}
Since $c_0,c_1\ge0$, this requires $\tilde e>1$, that is, $\Delta=b\,\tilde e>b$: silence sells above a verified efficient rating.

\emph{Step 3 (attention threshold, (iii)).} By~\eqref{eq:b_att}--\eqref{eq:premium_att}, $\tilde e=(C_0-C^*)/(C_0-C_1)$, so $\tilde e>1$ iff $C^*<C_1$. Substitute $C^*=C^{\rm def}+\psi(\bar C^{\rm sil}-C^{\rm def})$ and divide by $\bar C^{\rm sil}-C^{\rm def}>0$:
\begin{equation*}
    C^*<C_1\quad\Longleftrightarrow\quad\psi\;<\;\frac{C_1-C^{\rm def}}{\bar C^{\rm sil}-C^{\rm def}}.
\end{equation*}

\emph{Step 4 (the default must be optimistic, (iii)).} $C^*$ is a convex combination of $C^{\rm def}$ and $\bar C^{\rm sil}$, and $\bar C^{\rm sil}\ge C_1$ by~\eqref{eq:csil}. If $C^{\rm def}\ge C_1$, both points are at least $C_1$, and so is $C^*$. Hence $C^*<C_1$ requires $C^{\rm def}<C_1$.

\emph{Step 5 (full attention, (iv)).} With $\psi=1$,~\eqref{eq:etilde} gives $\tilde e=e$. Steps~1--2 are then the comparisons behind Proposition~\ref{prop:equilibrium}, with $b=\rho\omega_1\beta_1$ in the role of $\rho\beta_1$. Without hiding costs they are those of Proposition~\ref{prop:premium}.
